\section{Precipitation-type classification experiment}
\label{app:unet}

\subsection{Experiment design}\label{sec:4.8}

To test whether the correction helps a downstream precipitation product, we retrain an existing attention U-Net for precipitation type twice: once on uncorrected input (Model A) and once on parallax-corrected input (Model B). Only the input differs between the two.

\textbf{Data.} The U-Net works on 128 × 128 tiles of the INSAT-3DR L1C Mercator grid, with four channels: TIR1, TIR2, WV and MIR. Brightness temperatures are scaled as (BT − 180 K) / 150 K and clipped to {[}0, 1{]}. Each pixel carries the GPM DPR precipitation type of the footprint matched to it. Pixels without a GPM label are ignored in the loss and in the scores. GPM positions are on the ground, so parallax displaces the input but not the labels.

\textbf{Correction.} For Model B, all tiles of a scene are mosaicked, and the grid (regular in Mercator y and in longitude) is fitted from the tile coordinates to within \(10^{-4}\) pixel. Cloud-top height comes from TIR1 with the tropical standard atmosphere, with BT \textless{} 270 K taken as cloud. Every pixel is then corrected with the solver of Section~\ref{sec:3.3} and resampled onto the original pixel grid with the occlusion-aware method of Section~\ref{sec:3.5}, before the tiles are cut out again. Cells that receive no corrected pixel (0.06 \% of labelled pixels) are filled with the nearest corrected value for the input, and their labels are dropped for both models, so A and B are scored on the same pixels. The correction moves cloud by a median of 4.4 km, about one pixel, and changes TIR1 by more than 1 K on 34 \% of pixels.

\textbf{Splits.} The training set has 5,082 patches from 823 scenes (2018 and 2021--2023); the test set has 2,197 patches from 347 scenes (September 2023 and 2024). No scene, day or GPM overpass appears in both. There is no validation split: nothing is tuned, there is no early stopping, and the test set is read only after training has finished.

\textbf{Training.} Both models use the same network (about 1.1 million parameters, with layer normalisation), class-weighted cross-entropy, the Adam optimiser at a learning rate of \(10^{-6}\), a batch size of 1 and 100 epochs. They also start from the same initial weights and see the patches in the same order (seed 42). Each model is trained once. There are three tasks:

\begin{itemize}
\tightlist
\item
  four classes: stratiform, deep convective, shallow isolated and shallow non-isolated (class weights 1, 2, 3, 3);
\item
  stratiform versus convective (1, 2.2);
\item
  rain versus no rain (1, 3).
\end{itemize}

This gives six models.

\textbf{Scores.} Scores are computed per pixel on the test set: accuracy, macro F1, and per-class F1 and critical success index, CSI = TP / (TP + FP + FN). The difference B − A is tested with a paired bootstrap over test scenes (2,000 resamples), which keeps all patches of a scene together. The two-sided p-value is taken from the bootstrap distribution of the difference.

\subsection{Results}\label{sec:6.8}

Section~\ref{sec:6.4} leaves open whether correcting with the infrared height helps a precipitation product. Table~\ref{tab:4} compares the U-Net trained on uncorrected input (Model A) with the same network trained on parallax-corrected input (Model B), for the three tasks of Section~\ref{sec:4.8}.

\begin{longtable}[]{@{}
  >{\raggedright\arraybackslash}p{(\linewidth - 12\tabcolsep) * \real{0.1566}}
  >{\raggedright\arraybackslash}p{(\linewidth - 12\tabcolsep) * \real{0.2530}}
  >{\raggedleft\arraybackslash}p{(\linewidth - 12\tabcolsep) * \real{0.0602}}
  >{\raggedleft\arraybackslash}p{(\linewidth - 12\tabcolsep) * \real{0.0602}}
  >{\raggedleft\arraybackslash}p{(\linewidth - 12\tabcolsep) * \real{0.0964}}
  >{\raggedright\arraybackslash}p{(\linewidth - 12\tabcolsep) * \real{0.3133}}
  >{\raggedleft\arraybackslash}p{(\linewidth - 12\tabcolsep) * \real{0.0602}}@{}}
\caption{Test-set scores of Model A (uncorrected input) and Model B (parallax-corrected input): 2,197 patches from 347 scenes. B − A is given with its 95 \% interval from a paired bootstrap over scenes, and p is two-sided.}\label{tab:4}\\
\toprule\noalign{}
\begin{minipage}[b]{\linewidth}\raggedright
Task
\end{minipage} & \begin{minipage}[b]{\linewidth}\raggedright
Score
\end{minipage} & \begin{minipage}[b]{\linewidth}\raggedleft
A
\end{minipage} & \begin{minipage}[b]{\linewidth}\raggedleft
B
\end{minipage} & \begin{minipage}[b]{\linewidth}\raggedleft
B − A
\end{minipage} & \begin{minipage}[b]{\linewidth}\raggedright
95 \% interval
\end{minipage} & \begin{minipage}[b]{\linewidth}\raggedleft
p
\end{minipage} \\*
\midrule\noalign{}
\endfirsthead
\bottomrule\noalign{}
\endlastfoot
Four classes & Accuracy & 0.675 & 0.674 & −0.0004 & −0.0011 to +0.0004 & 0.290 \\*
& Macro F1 & 0.421 & 0.420 & −0.0012 & −0.0024 to 0.0000 & 0.050 \\*
& CSI stratiform & 0.700 & 0.701 & +0.0003 & −0.0004 to +0.0012 & 0.431 \\*
& CSI deep convective & 0.041 & 0.038 & −0.0030 & −0.0055 to −0.0007 & 0.008 \\*
& CSI shallow isolated & 0.265 & 0.264 & −0.0006 & −0.0019 to +0.0007 & 0.388 \\*
& CSI shallow non-isolated & 0.221 & 0.222 & +0.0010 & −0.0001 to +0.0020 & 0.070 \\*
Stratiform / convective & Accuracy & 0.709 & 0.709 & +0.0000 & −0.0012 to +0.0013 & 0.920 \\*
& Macro F1 & 0.686 & 0.685 & −0.0004 & −0.0014 to +0.0007 & 0.494 \\*
& CSI stratiform & 0.628 & 0.628 & +0.0006 & −0.0010 to +0.0022 & 0.481 \\*
& CSI convective & 0.429 & 0.428 & −0.0012 & −0.0023 to −0.0001 & 0.039 \\*
Rain / no rain & Accuracy & 0.863 & 0.863 & +0.0001 & −0.0002 to +0.0005 & 0.412 \\*
& Macro F1 & 0.718 & 0.717 & −0.0002 & −0.0007 to +0.0002 & 0.329 \\*
& CSI rain & 0.347 & 0.346 & −0.0005 & −0.0012 to +0.0001 & 0.120 \\*
\end{longtable}

Correcting the input does not improve the classifier. In all three tasks, accuracy and macro F1 of Model B are within 0.0012 of Model A. Only the convective classes change consistently, and they get slightly worse:

\begin{itemize}
\tightlist
\item
  deep convective CSI falls from 0.041 to 0.038 (B − A = −0.0030, 95 \% interval −0.0055 to −0.0007, p = 0.008);
\item
  convective CSI in the stratiform/convective task falls from 0.429 to 0.428 (−0.0012, −0.0023 to −0.0001, p = 0.039);
\item
  four-class macro F1 falls by 0.0012 (p = 0.05).
\end{itemize}

Stratiform scores and rain detection do not change significantly.

Three points qualify this result:

\begin{itemize}
\tightlist
\item
  \textbf{Size.} The changes are small. The largest relative change, 7 \% of deep convective CSI, is in a class that the network barely detects with either input (F1 0.08).
\item
  \textbf{Multiple tests.} With 22 scores tested, one or two p-values below 0.05 are expected by chance. None would pass a Bonferroni correction (p \textless{} 0.002).
\item
  \textbf{One run per model.} Each model was trained once, and run-to-run variation (for example from non-deterministic GPU operations) was not measured. It may be of similar size.
\end{itemize}

The result agrees with the only earlier comparison of this kind that we know of: a SEVIRI U-Net for rain rate, trained against GPM, also performed better without parallax correction \citep{dadderio2023}. Networks can learn a consistent displacement on their own \citep{hilburn2021}, so a geometric pre-correction adds little and brings its own errors (cloud-top height, holes and duplicates at occlusions).

The pattern is nevertheless consistent with Section~\ref{sec:5.5}. Uncorrected cold cores already sit on the GPM convective cells, so moving them by the full infrared parallax pulls the cold signal away from the convective labels, and convective scores drop slightly. Stratiform rain is horizontally extensive and insensitive to a one-pixel shift. For this classifier, correcting with the full infrared height brings no benefit. An effective height closer to the radar echo top (Section~\ref{sec:6.4}) can be tested with the same set-up.
